\documentclass[epjc3]{svjour3}

\usepackage{amsmath,amssymb,hyperref}
\usepackage{geometry}
\geometry{margin=1in}


\title{Geometric Origin of CP Phases from Hyperbolic Holonomy}
\author{Marvin L. Gentry}
\institute{Independent Researcher, Seattle, WA, USA \\
\email{drmlgentry@protonmail.com} \\
ORCID: 0009-0006-4550-2663}
\usepackage{amsmath,amssymb}
\usepackage{booktabs,graphicx,hyperref}
\journalname{Eur. Phys. J. C}

\begin{document}
\maketitle

\begin{abstract}
We show that CP-violating phases arise naturally as holonomy invariants
associated with flat $U(1)$ connections on compact hyperbolic $3$-manifolds.
If $M = \mathbb{H}^3/\Gamma$ admits an orientation-reversing isometry $\Theta$,
then $\Theta$ induces complex conjugation on $U(1)$ character representations
of $\pi_1(M)$. Any closed geodesic whose holonomy is not $\pm 1$
therefore provides a rephasing-invariant CP-odd observable.
For flat connections these phases are topological invariants and
are stable under all metric deformations preserving topological type.
We formulate the argument entirely in representation-theoretic language
and illustrate the mechanism using the Meyerhoff manifold ($\mathrm{OrientableClosedCensus}[1]$, $\mathrm{vol}=0.9814$, $H_1=\mathbb{Z}/5$).
\end{abstract}

The geometric construction developed here provides the theoretical
foundation for the companion results in~\cite{Gentry:CKM,Gentry:PMNS},
where the CKM and PMNS mixing matrices are derived explicitly from the
$K$- and $N$-factors of the Iwasawa decomposition of compact hyperbolic
3-manifold holonomy~\cite{Gentry:CKM,Gentry:PMNS}.

%--------------------------------------------------
\section{Hyperbolic Setup}
%--------------------------------------------------

Let $\Gamma$ be a torsion-free discrete group acting properly discontinuously
and cocompactly on $\mathbb{H}^3$ by orientation-preserving isometries.
Let
\[
M = \mathbb{H}^3/\Gamma
\]
be the resulting compact oriented hyperbolic $3$-manifold.

Closed geodesics correspond to conjugacy classes
$[\gamma] \subset \Gamma$ with $\gamma$ hyperbolic.

\begin{definition}[Flat $U(1)$ connection]
A flat $U(1)$ connection on $M$ is equivalently a character
\[
\rho : \pi_1(M) \longrightarrow U(1).
\]
The holonomy of a closed geodesic represented by $\gamma$
is
\[
W(\gamma) := \rho([\gamma]).
\]
\end{definition}

Because $U(1)$ is abelian, $W(\gamma)$ depends only on the
free homotopy class.

%--------------------------------------------------
\section{Orientation Reversal and CP}
%--------------------------------------------------

Let $\Theta : M \to M$ be an orientation-reversing isometry.
It induces an automorphism
\[
\Theta_* : \pi_1(M) \to \pi_1(M).
\]

\begin{definition}[CP action on characters]
Given a character $\rho$, define its CP-conjugate character by
\[
\rho^{\Theta}(\gamma)
=
\overline{\rho(\Theta_*^{-1}\gamma)}.
\]
\end{definition}

\begin{remark}
The inverse $\Theta_*^{-1}$ appears because $\Theta$ induces a
\emph{pullback} of representations: if $\Theta_*$ is an automorphism
of $\pi_1(M)$, the natural action on characters is by precomposition
with $\Theta_*^{-1}$, ensuring $\rho^\Theta$ is again a group homomorphism.
\end{remark}

\begin{lemma}[Holonomy conjugation]
Let $\rho$ be a flat $U(1)$ character.
For any $\gamma \in \pi_1(M)$,
\[
W_{\rho^{\Theta}}(\gamma)
=
\overline{W_{\rho}(\Theta_*^{-1}\gamma)}.
\]
\end{lemma}

\begin{proof}
By definition,
\[
W_{\rho^{\Theta}}(\gamma)
=
\rho^{\Theta}(\gamma)
=
\overline{\rho(\Theta_*^{-1}\gamma)}
=
\overline{W_{\rho}(\Theta_*^{-1}\gamma)}.
\]
\end{proof}

\begin{theorem}[Geometric CP violation]
Suppose:
\begin{enumerate}
\item $M$ admits an orientation-reversing isometry $\Theta$,
\item $\rho$ is a flat $U(1)$ character,
\item there exists a hyperbolic element $\gamma$ whose conjugacy class
is fixed by $\Theta_*$ (i.e., $\Theta_*[\gamma] = [\gamma]$) and such that
$W_\rho(\gamma) \neq \pm 1$.
\end{enumerate}
Then $\arg W_\rho(\gamma)$ is a rephasing-invariant
CP-odd observable.
\end{theorem}

\begin{proof}
Because the conjugacy class of $\gamma$ is fixed, we may choose a representative
so there exists $\delta\in\pi_1(M)$ with $\Theta_*^{-1}(\gamma)=\delta\gamma\delta^{-1}$.
Since $U(1)$ characters are conjugation-invariant,
$W_\rho(\Theta_*^{-1}\gamma) = W_\rho(\gamma)$ exactly.
Applying the CP action,
\[
W_\rho(\gamma) \;\xrightarrow{\text{CP}}\; W_{\rho^\Theta}(\gamma)
= \overline{W_\rho(\Theta_*^{-1}\gamma)} = \overline{W_\rho(\gamma)}.
\]
Hence $\arg W_\rho(\gamma) \mapsto -\arg W_\rho(\gamma)$.
If $W_\rho(\gamma) \neq \pm 1$, the phase is nonzero and flips sign under CP,
so it is a genuine CP-odd observable.  Rephasing invariance follows from
Proposition~\ref{prop:rephasing}.
\end{proof}

%--------------------------------------------------
\section{Rephasing Invariance}
\label{prop:rephasing}
%--------------------------------------------------

\begin{proposition}
Holonomy phases are invariant under all $U(1)$ gauge transformations.
\end{proposition}

\begin{proof}
Gauge transformations act by conjugation on $\rho$,
but since $U(1)$ is abelian, conjugation is trivial.
\end{proof}

\begin{remark}
The term \emph{rephasing invariance} here refers to $U(1)$ gauge
transformations of the flat connection, which are trivial for an abelian
structure group.  This is distinct from the sector-local rephasings
of the companion papers~\cite{Gentry:CKM,Gentry:PMNS}, which remove
unphysical phases from fermion field definitions.  Both notions confirm
that $\arg W_\rho(\gamma)$ is invariant and carries physical information.
\end{remark}

%--------------------------------------------------
\section{Topological Stability}
%--------------------------------------------------

\begin{theorem}
Let $\rho$ be a flat character.
Then $W_\rho(\gamma)$ depends only on the homotopy class of $\gamma$.
It is invariant under:
\begin{enumerate}
\item metric deformations preserving the topological type of $M$,
\item continuous deformations of $\rho$ within its connected component
in the character variety.
\end{enumerate}
\end{theorem}

\begin{proof}
Flatness implies holonomy depends only on the homotopy class.
Metric deformations do not change $\pi_1(M)$, so the holonomy of each class
remains unchanged.  Continuous variations of $\rho$ in character space
vary the holonomy continuously; if they stay within the same component,
the holonomy can change continuously but remains a topological invariant
in the sense that its value is determined by the representation.
For the purpose of CP violation, it suffices that the phase is
topologically quantised (e.g., a root of unity) and does not drift
under small metric perturbations.
\end{proof}

\begin{remark}[Discrete character variety]
When $H_1(M;\mathbb{Z})$ is purely torsion---as for the Meyerhoff
manifold with $H_1=\mathbb{Z}/5$---the character variety
$\mathrm{Hom}(\pi_1(M),U(1))$ consists of finitely many isolated points.
For $\mathbb{Z}/5$ there are exactly five, corresponding to the fifth
roots of unity.  Each connected component is a single point, so no
continuous deformation of $\rho$ is possible; the holonomy phase is
\emph{topologically quantized} and cannot drift.
\end{remark}

%--------------------------------------------------
\section{Example: The Meyerhoff Manifold}
%--------------------------------------------------

The Meyerhoff manifold $M=\mathrm{OrientableClosedCensus}[1]$
(SnapPy~\cite{SnapPy}, $\mathrm{vol}=0.9814$) is proven arithmetic
with invariant trace field $\mathbb{Q}(\sqrt{-283})$~\cite{Chinburg98}.
Its first homology group satisfies
\[
H_1(M;\mathbb{Z}) \cong \mathbb{Z}/5.
\]
The character variety consists of five isolated points; choose $\rho$
sending the generator to $e^{2\pi i/5}$ and let $\gamma$ represent the
generator of $\mathbb{Z}/5$.
Complex conjugation in $\mathbb{Q}(\sqrt{-283})$ induces an
orientation-reversing isometry acting by inversion
$k\mapsto -k$ on $H_1(M;\mathbb{Z})=\mathbb{Z}/5$,
mapping $\chi$ to its complex conjugate~\cite{MaclachlanReid,Chinburg98}.
The conjugacy class of $\gamma$ is fixed, so Theorem~2.3 applies.
Therefore
\[
W_\rho(\gamma) = e^{2\pi i/5},
\]
which is $\neq \pm 1$ and provides a concrete geometric CP-odd phase.

%--------------------------------------------------
\section{Conclusion}
%--------------------------------------------------

We have shown that CP-violating phases can arise from the geometry of
compact hyperbolic $3$-manifolds.  The holonomy of a flat $U(1)$ character
around a closed geodesic, when the manifold admits an orientation-reversing
isometry that fixes the geodesic’s conjugacy class, yields a rephasing-
invariant observable that flips sign under CP.  For flat connections these
phases are topological invariants, stable under metric deformations.
The Meyerhoff manifold provides an explicit example realising a fifth
root of unity as a CP-odd phase.
We note that the twist angles $\phi(\gamma)=\mathrm{Im}\log\lambda(\gamma)$
in the $A$-factor of the Iwasawa decomposition~\cite{Gentry:CKM,Gentry:PMNS}
are precisely the arguments of the flat $U(1)$ holonomies studied here;
the $\mathbb{Z}/5$ torsion quantizes these phases, providing the geometric
origin of CP violation in that framework.
Embedding this mechanism into a complete flavor model remains open.

\section*{Data Availability}
This manuscript contains no research data.

\section*{Conflict of Interest}
The author declares no conflicts of interest.

\begin{thebibliography}{99}

\bibitem{BrancoCP}
G.~C.~Branco, L.~Lavoura, and J.~P.~Silva,
\emph{CP Violation},
Oxford University Press, 1999.

\bibitem{BridsonHaefliger}
M.~R.~Bridson and A.~Haefliger,
\emph{Metric Spaces of Non-Positive Curvature},
Springer, 1999.

\bibitem{SnapPy}
M.~Culler, N.~Dunfield, M.~Goerner, and J.~R.~Weeks,
``SnapPy, a computer program for studying the geometry and topology
of 3‑manifolds,'' \url{http://snappy.computop.org}.

\bibitem{Chinburg98}
T.~Chinburg, E.~Friedman, K.~N.~Jones, and A.~W.~Reid,
``The arithmetic hyperbolic 3-manifold of smallest volume,''
\emph{Ann.\ Scuola Norm.\ Sup.\ Pisa} \textbf{30}, 1--40 (2001).

\bibitem{MaclachlanReid}
C.~Maclachlan and A.~W.~Reid,
\emph{The Arithmetic of Hyperbolic 3-Manifolds},
Springer, New York (2003).

\bibitem{Jarlskog}
C.~Jarlskog,
``A basis independent formulation of the connection between quark mass
matrices, CP violation and experiment,''
\emph{Z. Phys. C} \textbf{29} (1985), 491.

\bibitem{Ratcliffe}
J.~G.~Ratcliffe,
\emph{Foundations of Hyperbolic Manifolds},
Springer, 2006.

\bibitem{Thurston}
W.~P.~Thurston,
\emph{Three-Dimensional Geometry and Topology},
Princeton University Press, 1997.

\bibitem{Weeks}
J.~R.~Weeks,
``Computation of hyperbolic structures in knot theory,''
in \emph{Handbook of Knot Theory}, Elsevier, 2005.


\bibitem{Gentry:CKM}
M.~L. Gentry,
\newblock {CKM} quark mixing from geodesic axes of hyperbolic 3-manifold holonomy,
\newblock \emph{Phys. Rev. D} (2026), submitted March 2026, temporary {ID} EPJC-26-03-349.

\bibitem{Gentry:PMNS}
M.~L. Gentry,
\newblock Lepton mixing from Borel structure of hyperbolic holonomy,
\newblock \emph{Phys. Rev. D} (2026), submitted March 2026, temporary {ID} EPJC-26-03-351.
\end{thebibliography}

\end{document}